\chapter{Introduction}
\label{ch:introduction}

This chapter introduces the motivation, context, and objectives of this thesis, which focuses on the optimisation of virtual orbital spaces for post-Hartree-Fock methods and their implications for quantum computing. We begin by outlining the computational challenges in quantum chemistry that motivate the need for orbital optimisation. We then review existing strategies for virtual orbital reduction, leading to the OVOS method, which serves as the focus of this work. Finally, we state the specific research objectives and outline the thesis structure.

\section{Motivation}
\label{sec:int_motivation}

Chemistry fundamentally arises from the quantum mechanical interactions of electrons and nuclei \cite{SzaboOstlund1989}. The exact solution of the electronic Schrödinger equation for a molecule with $N$ electrons scales as $\mathcal{O}(N!)$ in the size of the configurational space, which rapidly renders exact calculations (such as Full Configuration Interaction) infeasible for any but the smallest systems \cite{Sherrill1999}. To achieve chemical accuracy (approximately 1 kcal/mol), it is essential to account for electron correlation effects beyond the mean-field Hartree-Fock (HF) approximation. Post-HF methods such as second-order Møller-Plesset perturbation theory (MP2), coupled cluster with single and double excitations (CCSD), and CCSD with perturbative triples [CCSD(T)] are widely used to capture these correlation effects \cite{Moller1934}. However, the computational cost of these methods scales steeply with system size, primarily due to the virtual orbital space. The number of virtual orbitals grows with the size of the basis set, and the computational cost of post-HF methods scales as $\mathcal{O}(O^2 V^2)$ for MP2 (AO to MO basis), $\mathcal{O}(V^4 O^2)$ for CCSD, and $\mathcal{O}(O^3 V^4)$ for CCSD(T), where $O$ is the number of occupied and $V$ is the number of virtual orbitals, notation from \cite{Helgaker2000}. As occupied orbitals are typically fewer than virtual orbitals, the virtual space dominates the computational cost. 

This challenge is not only relevant for classical computational chemistry but also has direct implications for quantum computing. Quantum algorithms for chemistry, such as the Variational Quantum Eigensolver (VQE) \cite{Peruzzo2014}, require a mapping of the electronic structure problem onto qubits \cite{Peruzzo2014}. The number of qubits required scales with the number of orbitals. Although current NISQ devices may have over 100 physical qubits (e.g., IBM's systems with 130+ qubits), device noise severely limits the number of qubits that can be reliably used. Reducing the number of orbitals (the virtual orbital space) is therefore crucial to minimise qubit requirements and mitigate the effects of noise, enabling chemistry simulations feasible in the near term \cite{Cao2019}. This motivates a closer examination of the virtual orbital space itself, which is the bottleneck for all post-HF methods and is the subject of the following section.

\section{The Virtual Orbital Space Problem}
\label{sec:int_virt_orb}

The virtual orbital space, which consists of the unoccupied molecular orbitals, plays a critical role in post-Hartree-Fock methods as it defines the space of excited configurations that contribute to electron correlation. The number of virtual orbitals increases with the size of the basis set, and their contributions to the correlation energy are not uniform. This section surveys various strategies developed to reduce or optimise the virtual orbital space, ultimately motivating the focus of this thesis on the Optimised Virtual Orbital Space (OVOS) method, originally proposed by Adamowicz and Bartlett in 1987 \cite{Adamowicz1987}.

The standard approach to defining virtual orbitals is to use the canonical Hartree-Fock (HF) virtual orbitals, which are obtained by diagonalising the Fock matrix \cite{Helgaker2000}. These orbitals are optimal for describing the single-determinant HF state. However, when we extend to multi-determinantal post-HF methods such as configuration interaction (CI) or coupled cluster (CC), the HF orbitals are no longer optimal for capturing electron correlation. This motivates the need to optimise the orbital space for correlated wavefunctions.

Natural orbitals, introduced by Löwdin in 1955 \cite{Lowdin1955}, are obtained by diagonalising the one-particle reduced density matrix (1-RDM) of an MP2 wavefunction. They are ordered by their occupation numbers, which can be used to truncate the virtual space. However, this approach requires an MP2 calculation to obtain the 1-RDM, creating a circular dependency.

Frozen natural orbitals (FNO) are a practical solution to this problem. They are computed from a low-cost correlated calculation, such as MP2, and the virtual orbitals are truncated based on an occupation threshold \cite{Taube2005}. This method is widely used in practice and can significantly reduce the computational cost while retaining accuracy.

The OVOS method takes a different approach by directly optimising a subset of the virtual orbitals to maximise the second-order correlation energy \cite{Adamowicz1987}. Unlike natural orbitals and FNO, OVOS does not necessarily require an initial MP2 calculation and can be systematically applied to any desired number of virtual orbitals. The OVOS method is the focus of this thesis, and we will implement and benchmark it in the context of unrestricted Hartree-Fock (UHF) calculations using the PySCF software package \cite{Sun2020}.

\section{Quantum Computing Context}
\label{sec:int_quantum_context}

Quantum computing has emerged as a promising paradigm for simulating quantum systems, including molecular electronic structure. Algorithms such as VQE can efficiently find ground state energies by preparing a parameterised quantum state and optimising the parameters using a classical optimiser \cite{Peruzzo2014}. However, current NISQ devices are severely limited in the number of qubits and the depth of quantum circuits they can reliably execute, primarily due to noise \cite{Cao2019}. This makes the connection between orbital-space reduction and quantum resource requirements directly relevant to this thesis.

In quantum chemistry, the electronic structure problem is typically mapped onto qubits using transformations such as the Jordan-Wigner \cite{Seeley2012}, which requires $Q = 2N_{\rm orb}$ qubits for $N_{\rm orb}$ spatial orbitals. Given the limited number of qubits and the need for shallow circuits, reducing the number of virtual orbitals directly translates to a reduction in qubit requirements and circuit depth \cite{Cao2019}. The OVOS method, by optimising a subset of the virtual orbitals, can significantly reduce the number of qubits needed for a given level of accuracy, $Q = 2N'_{\rm orb}$, where $N'_{\rm orb} < N_{\rm orb}$ is the number of orbitals retained after OVOS optimisation \cite{Adamowicz1987}.

\section{Objectives}
\label{sec:int_objectives}
The thesis will seek to achieve the following specific objectives:

\begin{enumerate}
    \item Implement and validate the OVOS algorithm of Adamowicz \& Bartlett \cite{Adamowicz1987} using the PySCF software package, for a representative set of molecules and basis sets.
    \item Characterise the energy recovery as a function of the active virtual orbital count for all test systems, quantifying the fraction of correlation energy recovered by the OVOS optimised orbitals.
    \item Investigate whether OVOS optimised molecular orbitals satisfy the Brueckner orbital condition at the MP2 level (vanishing T1 amplitudes), and assess the theoretical connection between OVOS and orbital-optimised MP2 \cite{Lee2018}.
    \item Quantify the potential benefit of OVOS optimised orbitals as a reference trial state for VQE type quantum computing algorithms: assess the expected improvement in reference-state overlap and compare initialisation strategies (default UHF, OVOS, UMP2 natural orbitals).
\end{enumerate}


\section{Thesis Outline}
\label{sec:int_outline}

The structure of this thesis follows. Chapter~\ref{ch:theory} reviews the theoretical background necessary to understand the OVOS method, including the relevance to quantum computing. Chapter~\ref{ch:method} describes the algorithmic implementation of the OVOS method, detailing code snippets of the full implementation. Chapter~\ref{ch:computational} details the computational setup, including the choice of test molecules, basis sets, and convergence criteria. Chapter~\ref{ch:results} presents the results of applying OVOS to a set of test molecules, analysing the energy recovery as a function of active virtual orbital count and using the OVOS orbitals for quantum computing  applications, particularly in the context of VQE algorithms. Finally, Chapter~\ref{ch:conclusion} summarises the key findings of the thesis, discusses the implications for quantum chemistry and quantum computing, and outlines directions for future research.


\section{Notation and Integral Conventions}
\label{sec:int_notation}

We establish the following specific notation and conventions for OVOS, maintaining consistency with the original Adamowicz \& Bartlett paper \cite{Adamowicz1987}.


\subsubsection{Index Ordering and Orbital Partitioning}


Orbital indices follows the standard convention:
\begin{itemize}
    \item $i,j,k,l,...$ for occupied orbitals
    \item $a,b,c,d,...$ for unoccupied orbitals
    \item $p,q,r,s,...$ for unrestricted orbitals
\end{itemize}

In OVOS, the virtual orbitals are further partitioned into active ($a,b,...$) and inactive ($e,f,...$) subsets, with the active space containing $N'$ orbitals optimised for correlation recovery.


\subsubsection{Integral Notation: Antisymmetrised vs Non-antisymmetrised}
% Transform (vir,occ|vir,occ) AO integrals to MO basis, assemble spin‑orbital
% and antisymmetrize. Returns <ab||ij> in physicist notation (a,b,i,j).

For integral notation, we distinguish between the antisymmetrised physicists' notation $\langle ab\|ij\rangle$ and the non-antisymmetrised chemists' notation $\langle ab|ij\rangle$. 
The antisymmetrised integral is defined as:
\begin{equation*}
\langle ab\|ij\rangle = \langle ab|ij\rangle - \langle ab|ji\rangle
\end{equation*}
where $\langle ab|ij\rangle$ is the standard two-electron integral in the MO basis. 
The physicists' notation is used throughout the OVOS derivations for consistency with the original literature and to simplify the expressions for the gradient and Hessian.


\subsubsection{Integral Notation: Chemists' vs.\ Physicists' Conventions}

The notation for two-electron integrals can differ between chemists and physicists, leading to potential confusion when comparing formulas from the literature with the output of quantum chemistry software. In chemists' notation, the integral is written as $(pq|rs)$, while in physicists' notation, it is written as $\langle pq | rs \rangle = (pr|qs)$ \cite{Helgaker2000}. The antisymmetrised integral, which is relevant for post-HF methods, is defined as:
\begin{equation*}
    \langle pq \| rs \rangle = \langle pq | rs \rangle - \langle pq | sr \rangle = (pr|qs) - (ps|qr)
\end{equation*}
The OVOS formulas in this thesis use the antisymmetrised physicists' notation, while PySCF returns integrals in the chemists' notation\cite{Sun2020}. The conversion between these notations is explicitly implemented in the code as: $\langle pq \| rs \rangle = (pr|qs) - (ps|qr)$. This establishes a clear and consistent framework for interpreting the integrals used in the OVOS method and ensures that the theoretical formulas align with the computational implementation.

